\documentclass[11pt,a4paper]{article}
\usepackage[utf8]{inputenc}
\usepackage[margin=1in]{geometry}
\usepackage{graphicx}
\usepackage{booktabs}
\usepackage{longtable}
\usepackage{hyperref}
\usepackage{natbib}
\usepackage{amsmath}
\usepackage{xcolor}

\title{Neurological Countermeasures for Spaceflight: A Transcriptomic-Driven, \\
       Diet-Based Neuroprotection Strategy}
\author{Biomni Analysis Pipeline}
\date{\today}

\begin{document}
\maketitle

\begin{abstract}
Spaceflight imposes unique stressors on the central nervous system, including microgravity, cosmic radiation, and circadian disruption. We analyzed rodent ISS brain RNA-seq data from 4 NASA OSDR studies (OSD-525, OSD-564, OSD-612, OSD-613) using random-effects meta-analysis to construct a consensus brain transcriptomic signature of \textbf{505 genes} (208 up, 297 down, 13 core). GSEA revealed dominant downregulation of mitochondrial oxidative phosphorylation (NES=-2.13), Complex I biogenesis (NES=-2.54), and protein homeostasis pathways, alongside upregulation of neuronal signaling (NES=+1.85), NMDA receptor activation (NES=+1.89), and long-term potentiation (NES=+1.66). LINCS L1000 connectivity mapping identified 221 candidate reverser compounds, including the nutrients sulforaphane (z=-5.89), ascorbic acid (z=-5.30), curcumin (z=-5.30), and luteolin (z=-4.31). Pharmacokinetic profiling of 7 brain-relevant nutrients identified ascorbic acid and sulforaphane as highest-priority countermeasures based on BBB penetration, brain tissue concentration, and spaceflight pathway relevance. We propose a phased vegan dietary countermeasure strategy (pre-flight preconditioning, in-flight protection, post-flight recovery) targeting mitochondrial dysfunction, oxidative stress, BBB disruption, and neuroinflammation. A critical finding is that curcumin, despite strong mechanistic evidence, has oral bioavailability ~1000-fold below therapeutic concentrations, requiring nanoformulation for CNS efficacy.
\end{abstract}

\section{Introduction}
Spaceflight induces neurological changes including hippocampal apoptosis, blood-brain barrier (BBB) disruption, and cognitive impairment \citep{Mao2020,Afshinnekoo2020}. Mitochondrial dysfunction has been identified as a central biological hub of spaceflight impact \citep{daSilveira2020}. We reanalyze rodent ISS brain transcriptomic data to characterize the neurological signature of spaceflight and identify evidence-based nutritional countermeasures.

\section{Methods}
\subsection{Data and Meta-Analysis}
Four rodent ISS brain RNA-seq datasets (OSD-525, OSD-564, OSD-612, OSD-613) were obtained from NASA OSDR. Random-effects meta-analysis was performed using \texttt{metafor} (REML estimator). Consensus genes required consistent direction across studies; core genes were significant in all 4 studies.

\subsection{Pathway Enrichment}
GSEA was performed with clusterProfiler against Hallmark and Reactome databases (p.adj $<$ 0.25). Pathways were classified into 10 neurological themes using keyword-based rules.

\subsection{LINCS L1000 Screening}
The consensus signature was queried against LINCS L1000 ($\sim$1.8M profiles). Negative connectivity z-scores indicate signature reversal. Compounds were annotated with Broad Drug Repurposing Hub data.

\subsection{Nutrient Pharmacokinetics}
PK data for 7 nutrients were compiled from published literature, including BBB penetration, oral bioavailability, peak plasma, brain concentration, and human dose.

\section{Results}
\subsection{Brain Consensus Signature}
The meta-analysis yielded 505 consensus genes (208 up, 297 down) with 13 core genes. The signature is skewed toward downregulation, reflecting suppression of metabolic and biosynthetic processes (Figure 1).

\subsection{Pathway Enrichment}
GSEA identified 23 significant Hallmark and 527 significant Reactome pathways. The dominant downregulated themes were mitochondrial dysfunction (mean NES=-1.57), protein homeostasis (mean NES=-1.57), and DNA damage repair (mean NES=-1.38). Upregulated themes included synaptic/neuronal signaling (mean NES=+1.53) and circadian/cholesterol dysregulation (mean NES=+1.73) (Figure 2, Table 1).

\subsection{LINCS Reversers}
LINCS screening identified 221 candidate reversers. Top drug reversers included vemurafenib (z=-9.68, RAF inhibitor) and chlorambucil (z=-7.24, DNA inhibitor). Nutrient reversers included sulforaphane (z=-5.89, Nrf2 activator), ascorbic acid (z=-5.30), curcumin (z=-5.30), and luteolin (z=-4.31) (Figure 3, Table 2).

\subsection{Nutrient Pharmacokinetics}
Comparative PK profiling (Figure 4, Table 3) identified ascorbic acid and sulforaphane as highest-priority countermeasures. Ascorbic acid is actively transported across the BBB via SVCT2, achieving $\sim$10-fold higher brain vs. plasma concentrations \citep{Hoffer2020}. Sulforaphane crosses the BBB within 15 minutes and activates Nrf2/ARE \citep{Zheng2022}. Curcumin has a critical bioavailability limitation: unconjugated plasma levels remain $\sim$1000-fold below in vitro effective concentrations \citep{Kroon2025}.

\subsection{Intervention Network}
A tripartite network (Figure 6) maps 7 nutrients to 5 target pathways and 4 protective mechanisms, revealing complementary multi-target coverage.

\subsection{Dietary Strategy}
Twelve curated vegan recipes (Figure 7, 8, Table 5) provide phased nutrient coverage across pre-flight, in-flight, and post-flight mission phases.

\section{Discussion}
The brain transcriptomic signature reveals a coherent pattern of mitochondrial energy failure, impaired protein homeostasis, and DNA repair suppression, consistent with the mitochondrial stress hub model \citep{daSilveira2020,Afshinnekoo2020}. The compensatory upregulation of neuronal signaling pathways may reflect adaptive plasticity but could predispose to excitotoxicity given impaired mitochondrial function.

The identification of sulforaphane and ascorbic acid as top nutrient reversers with favorable PK profiles provides a practical intervention strategy. Sulforaphane's Nrf2/ARE activation directly addresses the oxidative stress and mitochondrial dysfunction central to the signature, while ascorbic acid's active BBB transport ensures reliable brain accumulation.

The curcumin bioavailability limitation, highlighted by the 2025 critical reappraisal \citep{Kroon2025}, underscores the importance of PK-grounded nutrient selection. Mechanistic evidence alone is insufficient; brain-achievable concentrations must be confirmed.

\section{Limitations}
This analysis is based on rodent data; cross-species translation requires validation. PK data are from Earth-based studies and may not reflect spaceflight physiology. LINCS profiles cell line responses, not in vivo brain pharmacology. Dietary recommendations require validation in spaceflight analog studies.

\section{Conclusions}
Spaceflight induces a brain transcriptomic signature dominated by mitochondrial dysfunction and protein homeostasis failure. Seven vegan nutrients target these pathways through complementary mechanisms. Ascorbic acid and sulforaphane are the highest-priority countermeasures based on BBB penetration and brain concentration data. A phased dietary strategy provides practical implementation guidance for long-duration spaceflight.

\section*{Data Availability}
All scripts and data are available in the Zenodo package. Source data: NASA OSDR (OSD-525, OSD-564, OSD-612, OSD-613). LINCS L1000: NIH Common Fund. Pharmacokinetic data: published literature (see references).

\end{document}
